\documentclass[14pt,a4paper]{extarticle}

% --- XeLaTeX setup for Ukrainian ---
\usepackage{fontspec}
\usepackage{polyglossia}
\setmainlanguage{ukrainian}
\setotherlanguage{english}
\setmainfont{DejaVu Serif}
\setsansfont{DejaVu Sans}
\setmonofont{DejaVu Sans Mono}

% --- BibLaTeX setup ---
\usepackage[backend=biber, style=numeric, sorting=none]{biblatex}
\addbibresource{references.bib}

% --- Math ---
\usepackage{amsmath, amssymb, amsthm}
\usepackage{mathtools}
\usepackage{unicode-math}

% --- Layout ---
\usepackage[left=2.5cm, right=1.5cm, top=2cm, bottom=2cm]{geometry}
\usepackage{setspace}
\onehalfspacing

\usepackage{enumitem}
\usepackage{booktabs}
\usepackage{array}
\usepackage{hyperref}
\usepackage{titlesec}
\usepackage{xcolor}

\hypersetup{
    colorlinks=true,
    linkcolor=blue!60!black,
    citecolor=blue!60!black,
    urlcolor=blue!60!black
}

\titleformat{\section}{\large\bfseries}{\thesection.}{0.5em}{}
\titleformat{\subsection}{\normalsize\bfseries}{\thesubsection.}{0.5em}{}

% --- Theorem environments ---
\theoremstyle{definition}
\newtheorem{definition}{Визначення}[section]
\newtheorem{remark}{Зауваження}[section]

\theoremstyle{plain}
\newtheorem{theorem}{Теорема}[section]

% --- Custom commands ---
\newcommand{\vset}[1]{\mathcal{#1}}
\newcommand{\Real}{\mathbb{R}}

% ================================================================
\begin{document}

\begin{center}
    {\large\bfseries ДИПЛОМНА РОБОТА}\\[4pt]
    {\normalsize на тему:}\\[6pt]
    {\large\bfseries «Розробка та програмна реалізація алгоритму оптимізації\\
    розподілу вантажів у транспортній мережі»}\\[12pt]
\end{center}

\vspace{1.5cm}

% ================================================================
\section{Математична постановка задачі}\label{sec:formulation}

\subsection{Формалізація транспортної мережі}

Транспортну мережу формалізуємо як \textbf{мультимодальний зважений орієнтований граф} \cite{toth2014}:
\[
    G = (V,\; E,\; \Phi,\; \Psi),
\]
де:
\begin{itemize}[leftmargin=2em]
    \item $V = \{v_1, v_2, \ldots, v_n\}$ --- скінченна множина вузлів мережі;
    \item $E \subseteq V \times V$ --- множина орієнтованих ребер (доріг);
    \item $\Phi : V \to \vset{A}$ --- функція атрибутів вузлів;
    \item $\Psi : E \to \vset{B}$ --- функція атрибутів ребер.
\end{itemize}

\begin{definition}
\textbf{Вузол} $v_i \in V$ характеризується кортежем атрибутів:
\[
    \Phi(v_i) = \bigl(\mathrm{id}_i,\; \lambda_i,\; \mu_i,\; \delta_i^{\min},\; \delta_i^{\max},\; s_i\bigr),
\]
де $\mathrm{id}_i$ --- унікальний ідентифікатор; $(\lambda_i, \mu_i) \in \Real^2$ --- географічні координати (широта, довгота); $[\delta_i^{\min}, \delta_i^{\max}]$ --- часове вікно роботи хабу (наприклад, 08:00--21:00); $s_i \geq 0$ --- час обслуговування у вузлі (розвантаження, оформлення тощо).
\end{definition}

\begin{definition}
\textbf{Ребро} $e_{ij} = (v_i, v_j) \in E$ характеризується кортежем:
\[
    \Psi(e_{ij}) = \bigl(\ell_{ij},\; M_{ij}\bigr),
\]
де $\ell_{ij} \in \Real_{>0}$ --- довжина сегменту (метри); $M_{ij} \subseteq \mathcal{T}$ --- множина допустимих для цього ребра типів транспортних засобів.
\end{definition}

\subsection{Типи транспортних засобів та мультимодальність}

Визначимо множину типів транспортних засобів:
\[
    \mathcal{T} = \{\mathrm{Car},\; \mathrm{Bike},\; \mathrm{Walk}\},
\]
та визначимо підмножину ребер, доступних для кожного типу транспорту:
\[
    E_\theta = \bigl\{\, e_{ij} \in E \mid \theta \in M_{ij} \,\bigr\}, \qquad \theta \in \mathcal{T}.
\]
Підмножини $E_\theta$ можуть перетинатися: одне ребро може бути доступним для кількох типів транспорту.

Для кожного типу транспорту $\theta \in \mathcal{T}$ введемо характеристичну швидкість $\bar{u}_\theta$ (м/с) та вантажну місткість одиниці $Q_\theta$ (кг):
\[
    \bar{u}_{\mathrm{Car}} > \bar{u}_{\mathrm{Bike}} > \bar{u}_{\mathrm{Walk}}, \qquad
    Q_{\mathrm{Car}} \geq Q_{\mathrm{Bike}} \geq Q_{\mathrm{Walk}}.
\]

\textbf{Парк транспортних засобів.} Визначимо скінченну множину наявних транспортних засобів:
\[
    \mathcal{V} = \{w_1, w_2, \ldots, w_F\}, \qquad F = |\mathcal{V}|,
\]
де кожному транспортному засобу $w \in \mathcal{V}$ відповідає тип $\theta_w \in \mathcal{T}$. Кількість одиниць кожного типу:
\[
    F_\theta = \bigl|\{\, w \in \mathcal{V} \mid \theta_w = \theta \,\}\bigr|, \qquad \theta \in \mathcal{T}.
\]
Вантажна місткість транспортного засобу $w$ дорівнює $Q_{\theta_w}$.

\subsection{Формалізація задачі маршрутизації}

Нехай задано \textbf{депо} $v_0 \in V$ та множину пакунків $\mathcal{K} = \{k_1, k_2, \ldots, k_m\}$. Кожен пакунок $k \in \mathcal{K}$ характеризується парою:
\[
    k = \bigl(d_k,\; q_k\bigr),
\]
де $d_k \in V$ --- вузол призначення; $q_k \in \Real_{>0}$ --- маса пакунку. Усі пакунки відправляються з депо $v_0$.

Реальна адреса доставки може не збігатися з жодним вузлом графу $G$. У такому разі $d_k$ визначається як \textbf{найближчий вузол} графу до фактичних координат адреси:
\[
    d_k = \arg\min_{v \in V} \bigl\| (\lambda_v,\, \mu_v) - (\lambda_k^{\mathrm{addr}},\, \mu_k^{\mathrm{addr}}) \bigr\|,
\]
де $(\lambda_k^{\mathrm{addr}}, \mu_k^{\mathrm{addr}})$ --- геокодовані координати адреси. Залишкова відстань від $d_k$ до фактичної адреси долається пішки і враховується в механізмі «останньої милі» (\cite{wang2014}).

\subsubsection{Агрегований граф доставки}

Для формулювання задачі маршрутизації введемо \textbf{агрегований граф доставки} над множиною вузлів:
\[
    N = \{0,\, 1,\, 2,\, \ldots,\, m\},
\]
де вузол $0$ відповідає депо $v_0$, а вузол $i$ ($i = 1, \ldots, m$) --- пункту призначення пакунку $k_i$ (вузлу $d_{k_i} \in V$).

Для кожної пари вузлів $i, j \in N$ ($i \neq j$) та транспортного засобу $w \in \mathcal{V}$ визначимо \textbf{час переїзду} $t_{ij}^{(w)}$, що враховує механізм «останньої милі». Нехай $P_{ij} = (e_1, e_2, \ldots, e_{r_{ij}})$ --- найкоротший шлях від вузла $i$ до вузла $j$ у вуличному графі $G$ з $r_{ij}$ ребрами.

\textbf{Індекс розбиття} $s_{ij}^{(w)}$ визначається як максимальний префікс шляху, прохідний для типу $\theta_w$:
\[
    s_{ij}^{(w)} = \max\Bigl\{\, l \in \{0, \ldots, r_{ij}\} \;\Big|\; \forall\, l' \in \{1, \ldots, l\}:\; \theta_w \in M_{e_{l'}} \,\Bigr\}.
\]
Тобто транспортний засіб $w$ проїжджає ребра $e_1, \ldots, e_{s_{ij}^{(w)}}$ (основна частина), а ребра $e_{s_{ij}^{(w)}+1}, \ldots, e_{r_{ij}}$ долаються пішки («остання миля»). Зокрема, $s_{ij}^{(w)}=0$  означає, що жодне ребро не прохідне для $\theta_{w}$ і весь шлях долається пішки. Час переїзду:
\[
    t_{ij}^{(w)} = \underbrace{\sum_{l=1}^{s_{ij}^{(w)}} \frac{\ell_{e_l}}{\bar{u}_{\theta_w}}}_{\text{основна частина}} \;+\; \underbrace{\sum_{l=s_{ij}^{(w)}+1}^{r_{ij}} \frac{\ell_{e_l}}{\bar{u}_{\mathrm{Walk}}}}_{\text{«остання миля»}}.
\]

Позначимо \textbf{довжину «останньої милі»}:
\[
    \Delta_{ij}^{(w)} = \sum_{l=s_{ij}^{(w)}+1}^{r_{ij}} \ell_{e_l}.
\]

Шлях $P_{ij}$ вважається \textbf{здійсненним} для засобу $w$, якщо виконуються обидві умови:
\begin{enumerate}[leftmargin=2em, label=(\alph*)]
    \item кожне ребро «останньої милі» прохідне пішки: $\mathrm{Walk} \in M_{e_l}$ для $l > s_{ij}^{(w)}$;
    \item довжина «останньої милі» не перевищує порогу: $\Delta_{ij}^{(w)} \leq \Delta_{\max}$.
\end{enumerate}
Якщо хоча б одна з умов порушена, покладемо $t_{ij}^{(w)} = +\infty$. Параметр $\Delta_{\max}$ (типове значення --- 100~м) обмежує відстань, яку кур'єр може пронести вантаж пішки.

\subsubsection{Змінні рішення}

Введемо такі змінні:

\begin{enumerate}[leftmargin=2em]
    \item \textbf{Маршрутна змінна:}
    \[
        x_{ijw} \in \{0, 1\}, \quad i, j \in N,\; i \neq j,\; w \in \mathcal{V},
    \]
    де $x_{ijw} = 1$, якщо транспортний засіб $w$ переїжджає безпосередньо з вузла $i$ до вузла $j$.

    \item \textbf{Індикатор використання ТЗ:}
    \[
        y_w \in \{0, 1\}, \quad w \in \mathcal{V},
    \]
    де $y_w = 1$, якщо засіб $w$ задіяний у доставці.

    \item \textbf{Час прибуття:}
    \[
        \tau_i^{(w)} \in \Real_{\geq 0}, \quad i \in N,\; w \in \mathcal{V}
    \]
    --- момент часу прибуття засобу $w$ до вузла $i$ (допоміжна змінна для елімінації підциклів та контролю часових вікон).
\end{enumerate}

\subsection{Цільова функція та обмеження}

\textbf{Задача оптимізації} формулюється як задача змішаного цілочислового програмування:

\[
    \boxed{
    \min_{x,\, y} \; F
    = A \cdot \underbrace{\sum_{w \in \mathcal{V}} \sum_{\substack{i,j \in N \\ i \neq j}} t_{ij}^{(w)}\, x_{ijw}}_{T_{\mathrm{total}}}
    \;+\; B \cdot \underbrace{\sum_{w \in \mathcal{V}} y_w\vphantom{\sum_{\substack{i,j}}}}_{|\mathcal{V}_{\mathrm{act}}|}
    }
\]
де $A, B \geq 0$ --- вагові коефіцієнти ($A \gg B$): пріоритетом є мінімізація сумарного часу доставки $T_{\mathrm{total}}$, а вторинним критерієм --- мінімізація кількості задіяних ТЗ -- $|\mathcal{V}_{\mathrm{act}}|$.

\medskip

За умови виконання таких обмежень:

\begin{enumerate}[label=\textbf{О\arabic*.}, leftmargin=3em]
    \item \textbf{(Повнота доставки):} кожен пункт призначення відвідується рівно одним ТЗ:
    \[
        \sum_{w \in \mathcal{V}} \sum_{\substack{j \in N \\ j \neq i}} x_{ijw} = 1, \quad \forall\, i \in N \setminus \{0\}.
    \]

    \item \textbf{(Збереження потоку):} якщо ТЗ $w$ прибуває до вузла, він із нього і від'їжджає:
    \[
        \sum_{\substack{j \in N \\ j \neq i}} x_{jiw} = \sum_{\substack{j \in N \\ j \neq i}} x_{ijw}, \quad \forall\, i \in N,\; \forall\, w \in \mathcal{V}.
    \]

    \item \textbf{(Зв'язок із депо):} ТЗ виїжджає з депо тоді й лише тоді, коли він задіяний:
    \[
        \sum_{j \in N \setminus \{0\}} x_{0jw} = y_w, \qquad
        \sum_{j \in N \setminus \{0\}} x_{j0w} = y_w, \quad \forall\, w \in \mathcal{V}.
    \]

    \item \textbf{(Вантажна місткість):} сумарна маса пакунків у маршруті ТЗ $w$ не перевищує його місткість:
    \[
        \sum_{i \in N \setminus \{0\}} q_{k_i} \cdot z_{iw} \leq Q_{\theta_w}, \quad \forall\, w \in \mathcal{V},
    \]
    де $z_{iw} = \displaystyle\sum_{\substack{j \in N \\ j \neq i}} x_{ijw} \in \{0, 1\}$ --- індикатор того, що вузол $i$ обслуговується засобом $w$; $k_{i}$ --- пакунок, що відповідає вузлу $i \in N \setminus \{0\}$.

    \item \textbf{(Часові вікна та елімінація підциклів):} для кожного переїзду час прибуття узгоджується з часом проходження:
    \[
        \tau_j^{(w)} \geq \tau_i^{(w)} + s_i + t_{ij}^{(w)} - \mathcal{M}\bigl(1 - x_{ijw}\bigr), \quad \forall\, i,j \in N,\; i \neq j,\; \forall\, w \in \mathcal{V},
    \]
    де $s_i \geq 0$ --- час обслуговування у вузлі $i$ (визначений у кортежі $\Phi(v_i)$); $\mathcal{M} > 0$ --- достатньо велика константа. Час прибуття до кожного пункту має потрапляти у вікно роботи хабу:
    \[
        \delta_i^{\min} \leq \tau_i^{(w)} \leq \delta_i^{\max}, \quad \forall\, i \in N \setminus \{0\},\; \forall\, w \in \mathcal{V}.
    \]

    Це одночасно виконує дві функції: (1)~елімінація підциклів за схемою Міллера--Такера--Землін (MTZ) \cite{mtz1960}, оскільки змінні $\tau_i^{(w)}$ зростають вздовж маршруту; (2)~контроль часових вікон роботи хабів.

    \item \textbf{(Допустимість маршруту):} ТЗ $w$ може використати дугу $(i, j)$ лише якщо шлях у графі $G$ є здійсненним (основна частина прохідна для $\theta_w$, «остання миля» --- пішки):
    \[
        x_{ijw} = 1 \;\Rightarrow\; t_{ij}^{(w)} < +\infty, \quad \forall\, i,j \in N,\; i \neq j,\; \forall\, w \in \mathcal{V}.
    \]
\end{enumerate}

% ================================================================
\section{Метод розв'язання: генетичний алгоритм}\label{sec:ga}

Задача CVRPTW, сформульована в розділі~\ref{sec:formulation}, належить до класу NP-складних задач комбінаторної оптимізації~\cite{toth2014}. Для її наближеного розв'язання застосовано \textbf{генетичний алгоритм} (ГА) --- популяційну метаевристику, натхненну процесами природного відбору \cite{holland1975, goldberg1989}.

На відміну від конструктивних евристик (наприклад, послідовних вставок Соломона~\cite{solomon1987}), які будують єдиний розв'язок за один прохід, ГА одночасно підтримує сукупність (популяцію) кандидатів і поступово покращує її за допомогою операторів селекції, схрещування та мутації. Це дозволяє ефективно досліджувати простір розв'язків і знаходити якісні розв'язки для задач великої розмірності \cite{prins2004, braysy2005}.

\subsection{Подання розв'язку (хромосома)}

Використовується \textbf{перестановкове кодування} (giant tour): хромосома --- це впорядкована перестановка індексів усіх пунктів призначення без явних розмежувачів між маршрутами:
\[
    \pi = (\pi_1,\; \pi_2,\; \ldots,\; \pi_m), \qquad
    \pi_l \in \{1, 2, \ldots, m\}, \quad
    \pi_l \neq \pi_{l'} \text{ при } l \neq l'.
\]
Простір хромосом --- множина $\mathcal{S}_m$ усіх перестановок із $m$ елементів, $|\mathcal{S}_m| = m!$.

Розподіл пунктів між транспортними засобами здійснюється не на рівні кодування, а на рівні \textit{декодування}: процедура Split (див. наступний підрозділ) відновлює з хромосоми набір допустимих маршрутів.

\subsection{Декодер (процедура Split)}\label{subsec:split}

Перетворення хромосоми $\pi$ у набір маршрутів $\{R_1, \ldots, R_K\}$ виконує \textbf{жадібна процедура Split}. Алгоритм послідовно розміщує клієнтів у маршрути так, щоб кожен маршрут $R_r$ залишався допустимим для свого ТЗ $w_r$.

\textbf{Вхід:} хромосома $\pi$, парк $\mathcal{V} = (w_1, w_2, \ldots, w_F)$.

\textbf{Процедура:}
\begin{enumerate}[leftmargin=2em]
    \item Ініціалізувати: $r \leftarrow 1$, $R_r \leftarrow ()$, $\mathcal{U} \leftarrow \varnothing$ (множина нерозміщених клієнтів).
    \item Для кожного $l = 1, 2, \ldots, m$:
    \begin{enumerate}[label=(\alph*), leftmargin=2em]
        \item Побудувати кандидатний маршрут $R_r' = R_r \cdot (\pi_l)$ (додавання в кінець поточного маршруту).
        \item Якщо $R_r'$ допустимий для ТЗ $w_r$ (задовольняє обмеження О4 та О5 --- див. підрозділ~\ref{subsec:feasibility}), прийняти: $R_r \leftarrow R_r'$.
        \item Інакше закрити поточний маршрут (якщо $R_r \neq ()$, додати його до результату), перейти до наступного ТЗ: $r \leftarrow r+1$, $R_r \leftarrow ()$; повторити крок~(a).
        \item Якщо $r > F$ (парк вичерпано) і клієнта $\pi_l$ досі не розміщено, занести $\pi_l$ до $\mathcal{U}$.
    \end{enumerate}
    \item Якщо $R_r \neq ()$, додати до результату.
\end{enumerate}

\textbf{Вихід:} набір маршрутів $\{R_1, \ldots, R_K\}$ (де $K$ --- кількість задіяних ТЗ) та множина нерозміщених клієнтів $\mathcal{U}$.

\subsubsection{Перевірка допустимості маршруту}\label{subsec:feasibility}

Маршрут $R = (r_1, r_2, \ldots, r_k)$ ТЗ $w$ є \textbf{допустимим} тоді й лише тоді, коли виконуються наступні умови.

\textbf{(а) Вантажна місткість} (обмеження~О4):
\[
    \sum_{l=1}^{k} q_{k_{r_l}} \leq Q_{\theta_w}.
\]

\textbf{(б) Часові вікна} (обмеження~О5). Час прибуття $\tau_{r_l}$ та фактичний час початку обслуговування $b_{r_l}$ обчислюються рекурентно:
\[
    \tau_{r_l} = b_{r_{l-1}} + s_{r_{l-1}} + t_{r_{l-1},\, r_l}^{(w)},
    \qquad
    b_{r_l} = \max\bigl\{\tau_{r_l},\; \delta_{r_l}^{\min}\bigr\},
\]
де $r_0 = 0$ (депо), $b_{r_0} = \delta_0^{\min}$. Для всіх $l = 1, \ldots, k$ має виконуватись $\tau_{r_l} \leq \delta_{r_l}^{\max}$, а повернення до депо: $b_{r_k} + s_{r_k} + t_{r_k,\, 0}^{(w)} \leq \delta_0^{\max}$.

\textbf{(в) Допустимість у графі} (обмеження~О6): для кожної дуги $t_{r_{l-1}, r_l}^{(w)} < +\infty$ та $t_{r_k, 0}^{(w)} < +\infty$.

Якщо хоча б одна з умов (а)--(в) порушена, маршрут $R$ недопустимий.

\subsection{Фітнес-функція}

Якість хромосоми $\pi$ оцінюється \textbf{фітнес-функцією} (мінімізується):
\[
    \boxed{\;
    \Phi(\pi) = A \cdot T_{\mathrm{total}}(\pi) \;+\; B \cdot K(\pi) \;+\; P \cdot |\mathcal{U}(\pi)|
    \;}
\]
де:
\begin{itemize}[leftmargin=2em]
    \item $T_{\mathrm{total}}(\pi) = \sum_{r=1}^{K} \sum_{l=0}^{n_r} t_{r_l, r_{l+1}}^{(w_r)}$ --- сумарний час переїзду по всіх маршрутах;
    \item $K(\pi)$ --- кількість задіяних ТЗ;
    \item $\mathcal{U}(\pi)$ --- множина клієнтів, не розміщених у жодному маршруті;
    \item $A, B \geq 0$ --- вагові коефіцієнти цільової функції (див. розділ~\ref{sec:formulation});
    \item $P \gg A, B$ --- штрафний коефіцієнт за нерозміщеного клієнта.
\end{itemize}

Штрафний член $P \cdot |\mathcal{U}|$ забезпечує, що хромосоми з нерозміщеними клієнтами мають значно гірший фітнес за хромосоми з повним розміщенням. При достатньо великому $P$ (на практиці $P = 10^6$) будь-який повний розв'язок домінуватиме над неповним.

\subsection{Генетичні оператори}

\subsubsection{Селекція (турнірний відбір)}

Для вибору батьківських хромосом застосовується \textbf{турнірний відбір}: з популяції $\mathcal{P}$ випадково обирають $T$ особин, серед яких переможцем стає особина з найменшим фітнесом:
\[
    \mathrm{Tournament}(\mathcal{P}, T) = \arg\min_{\pi \in \mathcal{S}} \Phi(\pi),
    \qquad
    \mathcal{S} \subset \mathcal{P},\; |\mathcal{S}| = T.
\]
Параметр $T$ (розмір турніру) контролює селекційний тиск: більший $T$ прискорює збіжність, але зменшує різноманітність популяції.

\subsubsection{Схрещування (Order Crossover, OX)}

Оператор \textbf{OX}~\cite{oliver1987} створює нащадка з двох батьків $\pi^{(1)}, \pi^{(2)}$, зберігаючи відносний порядок генів. Алгоритм:

\begin{enumerate}[leftmargin=2em]
    \item Випадково обрати дві позиції $i, j$, $1 \leq i \leq j \leq m$.
    \item Скопіювати сегмент $\pi^{(1)}[i{:}j]$ у нащадка на ті самі позиції.
    \item Починаючи з позиції $(j \bmod m) + 1$, пройти $\pi^{(2)}$ у циклічному порядку. Кожен ген, що ще не присутній у нащадку, розмістити у наступній вільній позиції (також циклічно, починаючи з $(j \bmod m) + 1$).
\end{enumerate}

OX гарантує збереження перестановкової структури (усі гени унікальні) та успадковує просторові закономірності обох батьків.

\subsubsection{Мутація (swap)}

З імовірністю $p_m$ виконується \textbf{swap-мутація}: випадково обираються дві різні позиції $i, j \in \{1, \ldots, m\}$, значення яких міняються місцями:
\[
    (\pi_1, \ldots, \pi_i, \ldots, \pi_j, \ldots, \pi_m)
    \;\longrightarrow\;
    (\pi_1, \ldots, \pi_j, \ldots, \pi_i, \ldots, \pi_m).
\]
Цей оператор вносить локальні збурення та допомагає уникнути передчасного застрявання в локальних оптимумах.

\subsubsection{Елітизм}

Для запобігання втраті найкращих розв'язків при переході між поколіннями застосовується \textbf{елітизм}: $E$ найкращих хромосом поточної популяції копіюються в нову популяцію без змін. Це гарантує, що найменший фітнес у популяції не зростає між поколіннями (властивість монотонної збіжності найкращого значення).

\subsection{Ініціалізація популяції}\label{subsec:init}

Початкова популяція $\mathcal{P}_0$ розміру $\mathrm{POP}$ формується так:
\begin{itemize}[leftmargin=2em]
    \item \textbf{Одна хромосома} --- результат жадібного алгоритму найближчого сусіда (nearest-neighbour, NN): починаючи з депо, послідовно обирається ще не доданий клієнт із мінімальним часом переїзду від поточної точки.
    \item \textbf{Решта $\mathrm{POP} - 1$} --- випадкові перестановки множини $\{1, \ldots, m\}$.
\end{itemize}
Наявність NN-засіду прискорює ранню збіжність, не звужуючи різноманітності пошуку \cite{prins2004}.

\subsection{Основний алгоритм}

Генетичний алгоритм працює в парадигмі покоління-за-поколінням: на кожному кроці з поточної популяції $\mathcal{P}_g$ будується нова $\mathcal{P}_{g+1}$, поєднуючи елітизм та стохастичне породження нащадків.

\begin{enumerate}[leftmargin=2em, label=\textbf{Крок \arabic*.}]
    \item \textbf{Ініціалізація.} Побудувати $\mathcal{P}_0$ за правилами підрозділу~\ref{subsec:init}. Обчислити $\Phi(\pi)$ для кожного $\pi \in \mathcal{P}_0$. Запам'ятати найкращу особину $\pi^\star$.

    \item \textbf{Еволюційний цикл.} Для $g = 0, 1, \ldots, \mathrm{GEN}-1$:
    \begin{enumerate}[label=(\alph*), leftmargin=2em]
        \item \textbf{Елітизм.} Відсортувати $\mathcal{P}_g$ за зростанням $\Phi$ та скопіювати перших $E$ особин до $\mathcal{P}_{g+1}$.
        \item \textbf{Розмноження.} Поки $|\mathcal{P}_{g+1}| < \mathrm{POP}$:
        \begin{enumerate}[label=\roman*.]
            \item $\pi_a \leftarrow \mathrm{Tournament}(\mathcal{P}_g, T)$; \; $\pi_b \leftarrow \mathrm{Tournament}(\mathcal{P}_g, T)$.
            \item Якщо $\mathrm{rand}() < p_c$: $\pi' \leftarrow \mathrm{OX}(\pi_a, \pi_b)$;
                \quad інакше: $\pi' \leftarrow \pi_a$.
            \item Якщо $\mathrm{rand}() < p_m$: $\pi' \leftarrow \mathrm{Swap}(\pi')$.
            \item Додати $\pi'$ до $\mathcal{P}_{g+1}$.
        \end{enumerate}
        \item \textbf{Оцінювання.} Обчислити $\Phi(\pi)$ для всіх $\pi \in \mathcal{P}_{g+1}$.
        \item \textbf{Оновлення рекорду.} Якщо у $\mathcal{P}_{g+1}$ знайдено особину з меншим фітнесом, ніж $\Phi(\pi^\star)$, оновити $\pi^\star$.
    \end{enumerate}

    \item \textbf{Результат.} Застосувати декодер Split до $\pi^\star$, повернути набір маршрутів $\{R_1, \ldots, R_K\}$.
\end{enumerate}

\subsection{Параметри алгоритму}

ГА контролюється сімома параметрами (типові значення обрано згідно з рекомендаціями літератури~\cite{prins2004, goldberg1989}):

\begin{center}
\begin{tabular}{@{} l c p{7.5cm} @{}}
\toprule
\textbf{Параметр} & \textbf{Типове значення} & \textbf{Призначення} \\
\midrule
$\mathrm{POP}$ & $50$ & розмір популяції \\
$\mathrm{GEN}$ & $150$ & кількість поколінь \\
$p_c$ & $0{,}85$ & імовірність схрещування (OX) \\
$p_m$ & $0{,}15$ & імовірність мутації (swap) \\
$E$ & $2$ & розмір еліти \\
$T$ & $3$ & розмір турніру \\
$P$ & $10^{6}$ & штраф за нерозміщеного клієнта \\
\bottomrule
\end{tabular}
\end{center}

Вибір $p_c$ у діапазоні $[0{,}7; 0{,}9]$ і малого $p_m$ у $[0{,}05; 0{,}2]$ --- стандартна практика для GA на комбінаторних задачах: висока інтенсивність схрещування забезпечує ефективну рекомбінацію, а невелика мутація не дає популяції застрягти в локальному оптимумі, не руйнуючи знайдених структур.

\subsection{Обчислювальна складність}

Позначимо $m = |N \setminus \{0\}|$ --- кількість пунктів призначення.

\textbf{Одна оцінка фітнесу.} Декодер Split обробляє $m$ клієнтів. Для кожного клієнта перевірка допустимості маршруту (часові вікна, місткість) виконується інкрементно за $O(m)$ (повторне проходження поточного маршруту). Загалом: $O(m^2)$ на одне декодування + $O(m)$ на обчислення $T_{\mathrm{total}}$.

\textbf{Одне покоління.}
\begin{itemize}[leftmargin=2em]
    \item Породження $\mathrm{POP} - E$ нащадків: оператори OX і swap працюють за $O(m)$, тож усі нащадки --- $O(\mathrm{POP} \cdot m)$.
    \item Оцінювання нових хромосом: $O(\mathrm{POP} \cdot m^2)$.
\end{itemize}
Загалом одне покоління: $O(\mathrm{POP} \cdot m^2)$.

\textbf{Увесь ГА.} Загальна складність:
\[
    O\bigl(\mathrm{GEN} \cdot \mathrm{POP} \cdot m^2\bigr).
\]
За фіксованих $\mathrm{POP}$ і $\mathrm{GEN}$ (константи) асимптотична складність дорівнює $O(m^2)$ --- на порядок швидше за конструктивну евристику Соломона ($O(m^3)$).

У мультимодальному варіанті розрахунок матриці часів переїзду виконується окремо для кожного з $|\mathcal{T}|$ типів транспортних засобів. Оскільки $|\mathcal{T}|$ --- константа ($|\mathcal{T}| = 3$), асимптотика не змінюється.

Для задач із десятками та сотнями пунктів призначення час роботи становить секунди на сучасному обладнанні.

\begin{remark}
Генетичний алгоритм --- популяційна стохастична метаевристика, яка не гарантує знаходження глобального оптимуму, але на практиці забезпечує якісні субоптимальні розв'язки. На відміну від конструктивних евристик, ГА дозволяє повторно відвідувати простір розв'язків, виходити з локальних оптимумів завдяки мутації та ефективно рекомбінувати частини розв'язків через схрещування, що робить його придатним для задач CVRPTW середньої та великої розмірності \cite{prins2004}.
\end{remark}

\newpage
\printbibliography

\end{document}
